\subsection{Sprint 0}

Na preparação do ReadRace, participei das discussões iniciais com os demais AGES IV para compreender a proposta e organizar a condução do projeto. Em 13 de agosto, reunimo-nos para definir perguntas aos stakeholders e planejar o andamento da aula após a reunião. Em 15 de agosto, contribuí com a prototipação de telas e a escrita das user stories. Esse trabalho exigiu transformar a proposta ampla de incentivo à leitura em fluxos e decisões que pudessem orientar o desenvolvimento.

O planejamento documentado foi organizado em fundação do projeto, componentes compartilhados, biblioteca, detalhe e progresso de leitura, interação social, desafios e busca. Várias entregas dependiam de elementos comuns, como tema, componentes, schema e dados iniciais. O diagrama de dependências compartilhado entre os AGES IV ajudou a discutir a ordem de execução e os impactos de tarefas que permanecessem paradas.

Em 20 de agosto, mantive contato com os colegas de gestão e fiz uma leitura inicial da documentação relacionada às capas dos livros. Também participei da revisão da modelagem do banco. Em 25 de agosto, verifiquei discrepâncias entre as user stories da Sprint 1 e o modelo relacional, fiz anotações e conversei com os colegas para encaminhar os pontos aos AGES III. No dia seguinte, retomamos as dúvidas e discutimos mudanças, produzindo sugestões de melhoria. Essa experiência mostrou que requisitos e dados precisam evoluir de forma coordenada: uma regra pouco clara pode gerar interpretações distintas na API e na interface.

Outro eixo da minha participação foi a preparação das apresentações. Em 27 de agosto, trabalhei na criação dos slides para apresentar aos clientes os fluxos e as histórias priorizadas para a Sprint 1. A versão inicial ainda exigia ajustes, e sua construção me ajudou a organizar o que seria comunicado e a distinguir a visão completa do produto do recorte daquela entrega.

Essa preparação também teve importância pessoal. Apresentar para os stakeholders é um grande desafio para mim por conta da timidez. Nesta AGES, tenho participado ativamente das apresentações ao lado dos demais colegas de AGES IV, exercitando a exposição de ideias e a comunicação com pessoas externas à equipe. Considero esse contato uma oportunidade de desenvolvimento que vai além da implementação de software.
